\chapter{Literature Review \& Approach Selection (Khảo sát và chọn approach)}

\section{Cách khảo sát}

Khảo sát chia làm 2 cấp, vì đây là 2 loại kết quả khác nhau:
\begin{itemize}
  \item \textbf{Framework:} hệ thống đi từ dữ liệu thật đến mô phỏng, ví dụ NuRec, Cam2Sim, HUGSIM.
  \item \textbf{Method:} thuật toán cho một khối, ví dụ SplatAD hay DirectFisheye-GS cho khối dựng cảnh.
\end{itemize}

Một hệ thống được xem là framework Real2Sim đầy đủ khi có 5 phần: dữ liệu thật, tái dựng cảnh, road logic và kịch bản, tích hợp mô phỏng, và AD Stack chạy vòng kín.
Chỉ dựng được cảnh đẹp thì chưa đủ. Hệ thống phải có vòng mô phỏng phản ứng theo hành động của xe.
Tuy vậy, một framework nghiên cứu không nhất thiết có sẵn bước trích xuất bản đồ tự động hay tích hợp simulator ở mức sản phẩm.

Thông tin dưới đây được đối chiếu với paper, repository và tài liệu chính thức tháng 11/2026.
``Có'' nghĩa là tài liệu có nêu, không có nghĩa nhóm đã chạy thử trên dữ liệu của mình.
``Chưa xác nhận'' nghĩa là nguồn đã đọc không nêu, chưa kết luận là không có.

\section{Các framework Real2Sim}

Dưới đây là 6 framework chính có mã nguồn công khai.
Đây là danh sách chọn lọc, không phải toàn bộ các hệ thống hiện có.

\begin{table}[H]
  \centering
  \small
  \caption{Sáu framework Real2Sim chính}
  \begin{tabularx}{\linewidth}{>{\raggedright\arraybackslash}p{2.2cm} >{\raggedright\arraybackslash}p{2.1cm} >{\raggedright\arraybackslash}X >{\raggedright\arraybackslash}X}
    \toprule
    \textbf{Framework} & \textbf{Công bố} & \textbf{Điểm mạnh} & \textbf{Hạn chế} \\
    \midrule
    NuRec & NVIDIA, 2025--2026 & Tái dựng ở mức công nghiệp, có hệ sinh thái, nối được AlpaSim và CARLA & Tự nó không chạy mô phỏng, cần dữ liệu NCore và dữ liệu phụ trợ \\
    \addlinespace
    Cam2Sim & ASE 2026, Tools and Datasets & Từ log thật ra kịch bản CARLA, có OSM và OpenDRIVE & Chỉ có xe đỗ, chưa có traffic động. Mới thử 1 kịch bản, DAVE-2, 3 lần chạy mỗi miền \\
    \addlinespace
    HUGSIM & TPAMI 2026 & Simulator closed-loop, hơn 70 sequence, benchmark nhiều planner & Không dùng CARLA, HD map chỉ có ở nuScenes \\
    \addlinespace
    NeuroNCAP & ECCV 2024 & Kiểm thử an toàn theo kịch bản lấy cảm hứng từ Euro NCAP, dựng cảnh bằng NeRF & Tập trung đánh giá an toàn, không nhằm xây bản sao mạng đường \\
    \addlinespace
    WorldEngine & arXiv 2606.19836, 2026 & Gồm SimEngine (mô phỏng) và AlgEngine (huấn luyện), hướng post-training cho policy & Trọng tâm là huấn luyện policy, không phải so sánh chất lượng tái dựng \\
    \addlinespace
    AETHER & Repository 2026, chưa xác nhận paper & Module rõ ràng, traffic từ LimSim/SUMO hoặc IDM, nhiều planner & Dự án nhỏ, đang phát triển, cần tự kiểm tra khả năng tái lập \\
    \bottomrule
  \end{tabularx}
\end{table}

So với yêu cầu của đề bài, các framework đáp ứng như sau:

\begin{table}[H]
  \centering
  \footnotesize
  \caption{So sánh theo yêu cầu của project}
  \begin{tabularx}{\linewidth}{>{\raggedright\arraybackslash}p{1.85cm} >{\raggedright\arraybackslash}p{1.35cm} >{\raggedright\arraybackslash}p{1.6cm} >{\raggedright\arraybackslash}X >{\raggedright\arraybackslash}X >{\raggedright\arraybackslash}X >{\raggedright\arraybackslash}p{1.45cm}}
    \toprule
    \textbf{Framework} & \textbf{Fisheye} & \textbf{LiDAR} & \textbf{Bản đồ đường} & \textbf{Xe khác} & \textbf{Vòng kín} & \textbf{CARLA} \\
    \midrule
    NuRec & Có & Có & XODR tùy chọn trong USDZ & Có vật thể động & Qua AlpaSim; CARLA có điểm nối & Có \\
    \addlinespace
    Cam2Sim & Chưa xác nhận & Tùy chọn, dùng phát hiện xe đỗ & OSM, dùng \texttt{map.xodr} & Chỉ xe đỗ & DAVE-2 & Có \\
    \addlinespace
    HUGSIM & Chưa xác nhận & Không cần để dựng cảnh & Multi-plane ground, HD map ở nuScenes & Replay, IDM, hành vi tấn công & UniAD, VAD, LTF & Không \\
    \addlinespace
    NeuroNCAP & Chưa xác nhận & Chưa xác nhận & Chưa xác nhận & Theo kịch bản an toàn & End-to-end planner & Không \\
    \addlinespace
    WorldEngine & Chưa xác nhận & Chưa xác nhận & Chưa xác nhận & Mô hình hành vi sinh & Có & Không thấy nêu \\
    \addlinespace
    AETHER & Chưa xác nhận & Có & Bản đồ và mạng SUMO & LimSim/SUMO, IDM & UniAD, VAD, LTF, VLA & Không \\
    \bottomrule
  \end{tabularx}
\end{table}

Từ 2 bảng trên, có thể rút ra:
\begin{itemize}
  \item Không có framework nào tốt nhất ở mọi tiêu chí.
  \item Chỉ NuRec và Cam2Sim có tích hợp CARLA. Ví dụ NuRec + CARLA có điểm nhận lệnh lái nhưng chưa nối AD Stack (chương 11); vòng kín đầy đủ của NuRec được mô tả qua AlpaSim.
  \item Đề bài yêu cầu fisheye và LiDAR, nhưng hiện chỉ NuRec có tài liệu xác nhận cả hai.
  \item Traffic động là điểm yếu của Cam2Sim, nên không thể lấy Cam2Sim làm lời giải trọn vẹn cho yêu cầu traffic.
\end{itemize}

\section{Các hướng mới}

Các hệ thống sau chưa được tính vào 6 framework chính, nhưng hữu ích để tham khảo:

\begin{table}[H]
  \centering
  \small
  \caption{Các hướng mới đáng chú ý}
  \begin{tabularx}{\linewidth}{>{\raggedright\arraybackslash}p{2.3cm} l >{\raggedright\arraybackslash}X >{\raggedright\arraybackslash}X}
    \toprule
    \textbf{Hệ thống} & \textbf{Năm} & \textbf{Ý nghĩa} & \textbf{Vì sao chưa tính} \\
    \midrule
    Dash2Sim & 2026 & Từ video dashcam ra log 4D có tọa độ và tỉ lệ thật, chạy closed-loop trên nuPlan, đúng trường hợp đầu vào tối thiểu & Chưa kiểm tra mã nguồn; không dùng CARLA \\
    \addlinespace
    Instant NuRec & 2026 & Dựng cảnh feed-forward, khoảng 1,5 giây mỗi clip theo công bố của tác giả, tương thích AlpaSim & Thuộc hệ sinh thái NuRec \\
    \addlinespace
    Odyssey & 2026 & Benchmark closed-loop dài hơi: 100 kịch bản, mỗi kịch bản 100 giây từ nuPlan & Là benchmark; tác giả mới thông báo sẽ phát hành code \\
    \addlinespace
    UniSim & CVPR 2023 & Nền tảng về mô phỏng cảm biến neural closed-loop & Chưa tìm thấy mã nguồn công khai \\
    \bottomrule
  \end{tabularx}
\end{table}

\section{Ứng viên theo module}

Danh sách dưới đây là điểm bắt đầu, chưa phải lựa chọn cuối.

\begin{table}[H]
  \centering
  \caption{Ứng viên theo module}
  \begin{tabularx}{\linewidth}{l X}
    \toprule
    \textbf{Module} & \textbf{Các implementation có thể thử} \\
    \midrule
    Data Preparation & Dataset-specific converters, NuRec NCore adapters \\
    Pose Estimation & Ground-truth poses, LiDAR odometry (KISS-ICP, FAST-LIO), ViPE, Global SfM, COLMAP \\
    Reconstruction & NuRec, Instant NuRec, SplatAD, DirectFisheye-GS, Splatfacto, geometric reconstruction \\
    Rendering & gsplat, 3DGUT/3DGRUT, method-specific renderer \\
    Road Logic & OSM và OpenDRIVE, quỹ đạo kết hợp ảnh BEV (RoGS, RoadVGGT), existing maps \\
    Simulation & CARLA, AlpaSim, runtime riêng của HUGSIM hoặc AETHER \\
    Traffic & CARLA Traffic Manager, LimSim/SUMO, IDM \\
    AD Evaluation & CARLA controller, DAVE-2, UniAD, VAD, Autoware \\
    Metrics & Shared evaluation scripts \\
    \bottomrule
  \end{tabularx}
\end{table}

\section{Những vấn đề chưa được giải quyết trọn vẹn}

Ngay cả 6 framework trên cũng chưa giải quyết xong 4 vấn đề sau. Đây cũng là chỗ project có thể đóng góp:
\begin{itemize}
  \item \textbf{Tự động tái dựng road logic:} trích xuất đúng làn đường, kết nối, giao lộ, đèn tín hiệu và quyền ưu tiên từ dữ liệu thật.
  \item \textbf{Nhất quán hình học:} đưa cảnh tái dựng, quỹ đạo xe, LiDAR và bản đồ về cùng một hệ tọa độ có đơn vị thật.
  \item \textbf{Mô phỏng tất định:} giữ hành vi xe khác, thời gian mô phỏng và dữ liệu cảm biến lặp lại được để kiểm thử hồi quy.
  \item \textbf{Mở rộng quy mô:} giảm thời gian, chi phí tính toán, dung lượng và thao tác thủ công cho mỗi khu vực mới.
\end{itemize}

Khoảng trống chính là biến một cảnh tái dựng đẹp thành một môi trường mô phỏng tin cậy, tất định, đúng vị trí địa lý và điều khiển được.

\section{Tiêu chí chọn approach}

Bốn approach được chọn trong Phase 1 theo các tiêu chí sau, xếp từ quan trọng đến ít quan trọng:
\begin{enumerate}
  \item \textbf{Chạy được thật:} có mã nguồn, tài liệu, và chạy được với tài nguyên của nhóm.
  \item \textbf{Hợp dữ liệu đầu vào:} xử lý được pinhole, fisheye và LiDAR. Bộ 4 approach phải phủ đủ hai track.
  \item \textbf{Nối được với CARLA:} có sẵn hoặc có thể làm phần nối.
  \item \textbf{Hỗ trợ road logic và traffic:} có thể dùng lại cho yêu cầu mô phỏng.
  \item \textbf{Chi phí:} thời gian train, VRAM, lượng việc làm tay.
  \item \textbf{Chia được thành module:} có thể đưa vào tool và thay method.
\end{enumerate}

\section{Gợi ý ưu tiên}

Đây là gợi ý để nhóm thảo luận, chưa phải quyết định. Chưa có approach nào được nhóm chạy thử.

\begin{table}[H]
  \centering
  \small
  \caption{Gợi ý thứ tự nghiên cứu}
  \begin{tabularx}{\linewidth}{l l >{\raggedright\arraybackslash}X >{\raggedright\arraybackslash}X}
    \toprule
    \textbf{Thứ tự} & \textbf{Framework} & \textbf{Vì sao} & \textbf{Cần kiểm tra trước} \\
    \midrule
    1 & NuRec & Có tài liệu cho fisheye, LiDAR, CARLA & Yêu cầu GPU, chuyển dữ liệu sang NCore, giấy phép dữ liệu \\
    \addlinespace
    2 & Cam2Sim & Là tham khảo tốt cho việc nối bản đồ, tọa độ và CARLA & Chưa có traffic động; bước COLMAP làm thủ công \\
    \addlinespace
    3 & HUGSIM & Tham khảo tốt cho renderer, vật thể động, đánh giá closed-loop & Không dùng CARLA và LiDAR nên khó đưa vào pipeline chung \\
    \addlinespace
    4 & AETHER & Tham khảo tốt cho chia module và traffic & Độ ổn định, khả năng nhập scene mới \\
    \bottomrule
  \end{tabularx}
\end{table}

WorldEngine và NeuroNCAP phù hợp để tham khảo về kịch bản và đánh giá hơn là để đưa vào pipeline của project.
Các method riêng cho fisheye, như DirectFisheye-GS, vẫn cần được xem xét cho Track B.

\section{Nguồn}

{\raggedright
\begin{itemize}
  \item NuRec: \url{https://docs.nvidia.com/nurec/}, \url{https://github.com/NVIDIA/NuRec}
  \item NuRec + CARLA: \url{https://carla.readthedocs.io/en/latest/nvidia_nurec/}
  \item AlpaSim: \url{https://github.com/NVlabs/alpasim}
  \item Cam2Sim: \url{https://arxiv.org/abs/2607.04770}, \url{https://github.com/ast-fortiss-tum/cam2sim}
  \item HUGSIM: \url{https://doi.org/10.1109/TPAMI.2025.3647952}, \url{https://github.com/hyzhou404/HUGSIM}
  \item NeuroNCAP: \url{https://arxiv.org/abs/2404.07762}, \url{https://github.com/atonderski/neuro-ncap}
  \item WorldEngine: \url{https://arxiv.org/abs/2606.19836}, \url{https://github.com/OpenDriveLab/WorldEngine}
  \item AETHER: \url{https://github.com/siyingli99/Aether}
  \item Instant NuRec: \url{https://arxiv.org/abs/2607.14203}, \url{https://github.com/NVIDIA/instant-nurec}
  \item Odyssey: \url{https://arxiv.org/abs/2610.06469}
  \item UniSim: \url{https://waabi.ai/unisim/}
  \item Dash2Sim: \url{https://arxiv.org/abs/2606.07366}
  \item ViPE: \url{https://github.com/nv-tlabs/vipe}
  \item RoGS: \url{https://arxiv.org/abs/2405.14342}
  \item RoadVGGT: \url{https://arxiv.org/abs/2607.23758}
  \item OpenDRIVE Digital Twin Generator: \url{https://github.com/ftgTUGraz/opendrive-digital-twin-generator}
\end{itemize}
}

\section{Các approach được chọn và hypotheses}

\begin{table}[H]
  \centering
  \caption{Approach được chọn}
  \begin{tabularx}{\linewidth}{l l l X}
    \toprule
    \textbf{Approach} & \textbf{Tên} & \textbf{Track} & \textbf{Lý do chọn và hypothesis} \\
    \midrule
    A & TBD & TBD & TBD \\
    B & TBD & TBD & TBD \\
    C & TBD & TBD & TBD \\
    D & TBD & TBD & TBD \\
    \bottomrule
  \end{tabularx}
\end{table}
